% ===========================================================================
%  Annexe — Progression annuelle (édition du professeur uniquement)
%  Aucune date : une colonne est laissée libre pour les vôtres.
% ===========================================================================
\chapter{Progression annuelle}

Ce livre est rangé par domaine ; on n'enseigne pas dans cet ordre. Le tableau
qui suit propose la conversion en périodes, celles du calendrier scolaire.
Il ne porte aucune date : la colonne « vos dates » est là pour que vous
inscriviez les vôtres, année après année, sans que le livre se démode.

\section{Ce que dit le Programme d'Études, et ce que ce livre en fait}

Le Programme d'Études 2026 fixe, pour la classe de terminale L, à raison de
deux heures par semaine, trois composantes :

\begin{center}\footnotesize
\renewcommand{\arraystretch}{1.35}
\begin{tabular}{|L{3.4cm}|c|L{6.2cm}|}
\hline
\textbf{Composante} & \textbf{Durée} & \textbf{Contenus du PE}\\
\hline
Algèbre & 12 h & Second degré ; systèmes linéaires ($2\times2$ et
$3\times3$, méthode de Cramer)\\
Analyse & 30 h & Étude de fonctions (limites, dérivation, tracé) ;
logarithme népérien ; exponentielle ; primitives et intégration\\
Traitement des données et probabilités & 22 h & Statistique à deux
variables (ajustement affine, méthode de Mayer) ; probabilité (univers
discret fini)\\
\hline
\end{tabular}
\end{center}

Ce livre suit ce tableau à la lettre : trois parties, aucun chapitre au-delà.
Un point mérite d'être signalé. La Répartition Annuelle de l'année scolaire
en cours ajoute, elle, une période entière consacrée aux suites numériques,
absente du tableau des contenus du PE 2026 — c'est un choix éditorial
délibéré que ce livre ne les enseigne pas. Si votre établissement continue de
suivre la RAPE plutôt que le PE sur ce point précis, il faudra les enseigner
à partir d'un autre support.

\section{Progression proposée, en quatre périodes}

\begin{center}\footnotesize
\renewcommand{\arraystretch}{1.45}
\begin{tabular}{|c|L{5.9cm}|c|L{2.6cm}|}
\hline
\textbf{Période} & \textbf{Chapitres et contenus} & \textbf{Heures} &
\textbf{Vos dates}\\
\hline
\multirow{4}{*}{\textbf{1}}
 & Ch. 1 — Équations et inéquations du second degré \emph{(reprise rapide
   des acquis de première)} & 3 h & \\
 & Ch. 2 — Systèmes d'équations linéaires ($2\times2$ et $3\times3$,
   Cramer) & 5 h & \\
 & Ch. 3 — Lire et interpréter une courbe & 4 h & \\
 & \textbf{Devoir surveillé nº 1} et correction & 2 h & \\
\hline
\multirow{3}{*}{\textbf{2}}
 & Ch. 4 — Limites et branches infinies & 5 h & \\
 & Ch. 5 — Dérivation et sens de variation & 5 h & \\
 & Ch. 6 — Étude complète d'une fonction & 5 h & \\
\hline
\multirow{4}{*}{\textbf{3}}
 & Ch. 7 — La fonction logarithme népérien & 5 h & \\
 & Ch. 8 — La fonction exponentielle & 4 h & \\
 & Ch. 9 — Primitives et calcul intégral & 2 h & \\
 & \textbf{Devoir surveillé nº 2} et correction & 3 h & \\
\hline
\multirow{4}{*}{\textbf{4}}
 & Ch. 10 — Dénombrement & 4 h & \\
 & Ch. 11 — Probabilité & 6 h & \\
 & Ch. 12 — Statistique à deux variables & 4 h & \\
 & \textbf{Devoir surveillé nº 3} — sujet complet, en baccalauréat blanc &
   2 h & \\
\hline
\end{tabular}
\end{center}

\noindent Total : environ cinquante-neuf heures, marges de révision
comprises. Les volumes du Programme d'Études — trente heures d'analyse,
douze d'algèbre, vingt-deux de traitement des données — y sont respectés.
Quatre périodes suffisent au contenu du programme ; la cinquième période du
calendrier scolaire, que la RAPE réserve aux suites numériques, peut servir
de période de révision générale et de baccalauréat blanc.

\section{Repères de fin de période}

À la fin de chaque période, un élève qui suit doit savoir faire, seul, ce qui
suit. Si ce n'est pas le cas, c'est ici qu'il faut revenir, avant d'avancer.

\begin{center}\footnotesize
\renewcommand{\arraystretch}{1.4}
\begin{tabular}{|c|L{10.4cm}|}
\hline
\textbf{Fin de} & \textbf{L'élève sait}\\
\hline
P1 & résoudre une équation et une inéquation du second degré ; résoudre un
     système $3\times3$ par Cramer ; lire sur une courbe un ensemble de
     définition, un extremum, un signe, et dresser un tableau de variation à
     partir d'un dessin\\
P2 & calculer une limite, nommer une asymptote, dériver un quotient,
     étudier le signe de la dérivée et mener une étude complète jusqu'au
     tracé\\
P3 & résoudre une équation en $\ln$ ou en $e^{x}$ ; lever l'indétermination
     par la factorisation en $x$ ; vérifier qu'une fonction est une
     primitive et calculer une aire\\
P4 & choisir entre $n^{p}$, $\arrang{p}{n}$ et $\combi{p}{n}$ ; calculer une
     probabilité et passer par l'événement contraire ; tracer un nuage à
     l'échelle imposée et déterminer la droite de Mayer\\
\hline
\end{tabular}
\end{center}

\section{Points de vigilance, période par période}

\textbf{Période 1.} Le tableau de signes reste, chaque année, le point
faible des copies — bien avant le calcul du discriminant. La règle de Sarrus
demande une séance entière pour être posée sans erreur de signe ; faites
recopier les deux colonnes supplémentaires au crayon.

\textbf{Période 2.} Les élèves savent dériver longtemps avant de savoir
conclure. Exigez la phrase « le dénominateur étant un carré, il est
strictement positif ; le signe de $f'$ est donc celui de son numérateur »,
qui figure telle quelle dans les corrigés officiels. Exigez également le
tracé sur papier millimétré, à l'unité imposée.

\textbf{Période 3.} Deux automatismes valent la moitié des points du
problème : la factorisation par $x$ qui fait apparaître $\frac{\ln x}{x}$,
et les conditions d'existence écrites avant toute résolution.

\textbf{Période 4.} La difficulté n'est jamais le calcul, toujours la
lecture de l'énoncé : « simultanément », « successivement avec remise »,
« successivement sans remise ». Faites entourer ces mots avant tout calcul.
Pour la statistique, la rupture d'échelle sur l'axe des ordonnées vaut un
point entier.

\section{Si le temps manque}

Il manquera. Voici, par ordre de priorité décroissante, ce qu'il faut
préserver quoi qu'il arrive, au regard du barème de l'épreuve :

\begin{enumerate}[label=\textbf{\arabic*.}, leftmargin=*, itemsep=0.2em]
  \item \textbf{Le chapitre 11} (probabilité) et le chapitre 10 qui le
        prépare : cinq points une session sur deux.
  \item \textbf{Le chapitre 12} (statistique) : cinq points l'autre session
        sur deux, et le plus accessible de tous.
  \item \textbf{Les chapitres 5 et 7} (dérivation, logarithme) : le c\oe ur
        du problème, six à sept points sur dix.
  \item \textbf{Le chapitre 9} : deux points et demi, pour deux heures
        d'enseignement seulement — le meilleur rendement horaire du
        programme.
\end{enumerate}

Les chapitres 3 et 4, s'ils doivent être abrégés, peuvent l'être en
travaillant directement sur les figures du chapitre 6 : les lectures
graphiques s'y font en situation, et le temps gagné est considérable.
